% ========================================
% 共通プリアンブル（講義MODE・演習MODE）
% パッケージ／解答切り替え／番号／空欄／定義・公式・定理の囲み
% ========================================
\usepackage{amsmath, amssymb, mathtools, bm}
\usepackage{ascmac}
\usepackage[dvipdfmx]{graphicx}
\usepackage[dvipdfmx]{xcolor}
\usepackage{tikz}
\usetikzlibrary{calc, arrows.meta}
\usepackage[most]{tcolorbox}
\usepackage{enumitem}
\usepackage{tasks}

\linespread{1.25}
\everymath{\displaystyle}
\setlength{\parindent}{0pt}
\setlist[enumerate,1]{label=（\arabic*）, leftmargin=*, labelsep=0.5em}

% ========================================
% 解答表示の切り替え（既定は教師用．make student で \studentmode が定義され生徒用になる）
% ========================================
\newif\ifshowanswer
\showanswertrue
\ifdefined\studentmode\showanswerfalse\fi

% ========================================
% 番号（「章.節.番号」形式，節ごとにリセット）
% ========================================
% jsarticle には章がないので，章番号は専用カウンタで持つ
\newcounter{mathchapter}
\newcounter{definition}[section]
\newcounter{formula}[section]
\newcounter{theorem}[section]
\renewcommand{\thedefinition}{\arabic{mathchapter}.\arabic{section}.\arabic{definition}}
\renewcommand{\theformula}{\arabic{mathchapter}.\arabic{section}.\arabic{formula}}
\renewcommand{\thetheorem}{\arabic{mathchapter}.\arabic{section}.\arabic{theorem}}

% ========================================
% 書き込み用コマンド
% ========================================
% 解答は青字・太字（数式も和文も太字になるよう \bfseries\boldmath を使う）
% 色は外側から \textcolor で掛けると数式中の \text{} の和文に効かないため，
% 箱の内側で \color を指定する
\newcommand{\answerstyle}[1]{{\color{blue}\bfseries\boldmath #1}}

% 空欄の左右それぞれの余白
\newcommand{\fitblankpad}{2zw}

% 空欄（答えの幅＋左右の余白で下線を引く）
% 使い方: \fitblankbf[]{答え}  ※文字列は \text{...} で囲む
\newsavebox{\fitblanksavebox}
\newcommand{\fitblankbf}[2][]{%
  \sbox{\fitblanksavebox}{\answerstyle{$\displaystyle #2$}}%
  \underline{\makebox[\dimexpr\wd\fitblanksavebox+\fitblankpad*2\relax][c]{%
    \ifshowanswer\usebox{\fitblanksavebox}%
    \else\phantom{\usebox{\fitblanksavebox}}\fi}}%
}

% 下線なしの空欄（分数の分子・分母など，下線と分数線が重なる場所で使う）
% 使い方: \frac{\fitblankbox{y \text{ の変化量}}}{\fitblankbox{x \text{ の変化量}}}
\newcommand{\fitblankbox}[1]{%
  \sbox{\fitblanksavebox}{\answerstyle{$\displaystyle #1$}}%
  \makebox[\dimexpr\wd\fitblanksavebox+\fitblankpad*2\relax][c]{%
    \ifshowanswer\usebox{\fitblanksavebox}%
    \else\phantom{\usebox{\fitblanksavebox}}\fi}%
}

% 数式中の計算欄（下線なし，答えの幅＋左右2em を確保）
% 使い方: $= \answermath{3 - 1 = 2}$
\newcommand{\answermath}[1]{%
  \ifshowanswer\mbox{\answerstyle{$\displaystyle #1$}}%
  \else\hspace{2em}\phantom{#1}\hspace{2em}\fi
}

% 記述・計算スペース（薄い枠で高さを確保）
% 使い方: \writespace[高さ]{解答例}
\newcommand{\writespace}[2][5em]{%
  \begin{tcolorbox}[enhanced, colback=white, colframe=black!35,
      boxrule=0.4pt, sharp corners, height=#1, valign=center,
      left=8pt, right=8pt]
    \ifshowanswer\answerstyle{#2}\fi
  \end{tcolorbox}%
}

% ========================================
% 定義（赤系）・公式（青系）・定理（青系・太枠）
% 使い方: \begin{definitionbox}{用語} ... \end{definitionbox}
% ========================================
\tcbset{
  mathbox/.style={
    enhanced, sharp corners,
    left=10pt, right=10pt, top=8pt, bottom=8pt,
    fonttitle=\bfseries,
  }
}
\newtcolorbox{definitionbox}[1]{
  mathbox, colback=red!5, colframe=red!70!black, boxrule=0.8pt,
  code={\refstepcounter{definition}},
  title={定義 \thedefinition\quad #1}
}
\newtcolorbox{formulabox}[1]{
  mathbox, colback=blue!5, colframe=blue!70!black, boxrule=0.8pt,
  code={\refstepcounter{formula}},
  title={公式 \theformula\quad #1}
}
\newtcolorbox{theorembox}[1]{
  mathbox, colback=blue!5, colframe=blue!70!black, boxrule=1.2pt,
  code={\refstepcounter{theorem}},
  title={定理 \thetheorem\quad #1}
}
% テクニック（緑系）
% 使い方: \begin{techniquebox}[ラベル]{見出し} ... \end{techniquebox}
\newcounter{technique}[section]
\renewcommand{\thetechnique}{\arabic{mathchapter}.\arabic{section}.\arabic{technique}}
\newtcolorbox{techniquebox}[2][]{
  mathbox, colback=green!5, colframe=green!70!black, boxrule=0.8pt,
  code={\refstepcounter{technique}\if\relax\detokenize{#1}\relax\else\label{#1}\fi},
  title={Technique \thetechnique\if\relax\detokenize{#2}\relax\else\quad\textbf{#2}\fi}
}
